% UNIFIED 2026-05-13
% =====================================================================
% Manuscript section draft -- Agent 3
% Preliminaries: Torsion, rearrangement, and deficits
%
% This file is to be merged by Agent 17 into
% Manuscript/01_preliminaries_draft.tex.
%
% Canonical notation: Manuscript/00_notation_register.md.
% Standing convention here:
%   - dimension is $n\ge 2$ (not $N$);
%   - the BDV-normalised energy is $E(\Omega)=-T(\Omega)/2$;
%   - the Saint--Venant deficit is $\delta_T(\Omega)=E(\Omega)-E(B^*)\ge 0$;
%   - the Faber--Krahn deficit is
%     $\delta_{\mathrm{FK}}(\Omega)=(|\Omega|/\omega_n)^{2/n}(\lambda_1(\Omega)-\lambda_1(B^*))$;
%   - the Fraenkel asymmetry is $\Fr(\Omega)\in[0,2)$;
%   - the foliation is $E_\rho=\{u>t(\rho)\}$ with $|E_\rho|=\omega_n\rho^n$.
%
% =====================================================================
% SOURCE MAP
% Item                                              Status     Pointer
% ----------------------------------------------------------------------
% B1: torsion energy E, rigidity T, E=-T/2          PROVED HERE
%       (definitions and equivalence), with         CITE       Brasco--De Philippis--
%       sign convention fixed.                                 Velichkov 2015 [BDV],
%                                                              Final/SaintVenantAssembly.tex,
%                                                              Final-BoundedReduction,
%                                                              Final-KohlerJobinTransfer
%                                                              for the two coexisting
%                                                              conventions.
%
% B2: Saint--Venant deficit delta_T,                ADAPTED    Final-BoundedReduction
%       scale-invariant form D,                                §1; WIP_master.tex §II.
%       Saint--Venant inequality.                  CITE        Polya 1948; Talenti 1976.
%
% B3: Fraenkel asymmetry A(Omega), [0,2)            CITED      Maggi 2008 survey,
%                                                              Final/SaintVenantAssembly.tex.
%
% B4: Faber--Krahn deficit delta_FK,                PROVED HERE (normalisation lemma);
%       eigenvalue normalisation                    CITED      Henrot 2006; KJT.tex.
%
% B5: Talenti pointwise comparison                  CITED       Talenti 1976, Thm.1.
%
% B6: volume-radius foliation E_rho = {u>t(rho)},   PROVED HERE properties of t(rho)
%       monotonicity of t                                       follow from coarea;
%                                                  ADAPTED     WIP_LevelSetIdentity.tex
%                                                              §1, WIP_master.tex §II.
%
% B7: profile-gap estimate (sup-norm Talenti bound)  CITED       Talenti 1976 Thm 1.
%
% Profile-gap estimate used by Plan 2 (h(s)=b(s)-v(s)
%   bound and consequence on t(rho))                PROVED HERE Plan 2/level-set-deficit-identity.md §4, Lem 8.1
%
% Bad-set Lebesgue estimate (set where -t'(rho)<tau_0
%   has measure <= K delta_T)                       PROVED HERE WIP_WeightedMetricTrace.tex
%                                                              Lemma "Bad-set" §4
%                                                              ("integrated Talenti gap").
%
% B9: Kohler--Jobin inequality
%       T(Omega) lambda_1(Omega)^{(n+2)/2}
%         >= T(B^*) lambda_1(B^*)^{(n+2)/2}        CITED       Brasco 2014, Thm 3.3;
%                                                              orig. Kohler--Jobin 1978;
%                                                              Final-KohlerJobinTransfer
%                                                              Thm. 2.1.
%
% Elementary one-variable inequality              PROVED HERE  Final-KohlerJobinTransfer
%   (1-s)^{-p}-1 >= ps                                          Lem. 4.1 (reproduced).
% =====================================================================

\subsection{Preliminaries II: Torsion, rearrangement, and deficits}
\label{sec:prelim-torsion}

This section fixes the torsion and Faber--Krahn functionals used
throughout, the Saint--Venant and Faber--Krahn deficits, the Fraenkel
asymmetry, the Talenti comparison, the volume--radius foliation
$\Erho=\{u>t(\rho)\}$, the profile-gap and bad-set bounds used in the
Plan~2 weighted metric trace, and the Kohler--Jobin inequality used in
the final transfer.

Conventions follow \texttt{00\_notation\_register.md}: we write the
ambient dimension as~$n\ge 2$, the BDV-normalised Saint--Venant energy
as $E(\Om)=-T(\Om)/2$, the Saint--Venant deficit as
$\deltaT(\Om)=E(\Om)-E(\Bstar)\ge 0$, the Faber--Krahn deficit as
$\deltaFK(\Om)=(|\Om|/\omega_n)^{2/n}(\lambda_1(\Om)-\lambda_1(\Bstar))$,
and the Fraenkel asymmetry as $\Fr(\Om)\in[0,2)$.

\subsubsection{Torsion potential, rigidity, and energy}
\label{ssec:torsion}

Throughout, $\Om\subset\R^n$ is open with $0<|\Om|<\infty$. No
connectedness, smoothness, or boundedness is assumed unless stated.

\begin{definition}[Torsion potential]\label{def:utorsion}
The \emph{variational torsion function} of $\Om$ is the unique
minimiser
\[
   u_\Om\in H^{1}_{0}(\Om),\qquad
   u_\Om=\arg\min\Bigl\{
   J(v)\,:\,v\in H^{1}_{0}(\Om)
   \Bigr\},
   \qquad
   J(v):=\tfrac12\!\int_{\Om}\!|\nabla v|^{2}\,dx-\!\int_{\Om}\!v\,dx.
\]
Equivalently, $u_\Om$ is the weak solution of $-\Delta u_\Om=1$ in
$\Om$ with $u_\Om=0$ on $\partial\Om$ in the $H^{1}_{0}$ trace sense.
Testing the weak equation against $u_\Om^{-}$ yields $u_\Om\ge 0$
a.e.\ in $\Om$.
\end{definition}

\begin{remark}[Existence]
The functional $J$ is strictly convex, coercive on $H^{1}_{0}(\Om)$ by
the Poincar\'e--Sobolev inequality
$\|v\|_{L^{2^*}}\le C(n)\|\nabla v\|_{L^{2}}$ on sets of finite
measure (valid since $|\Om|<\infty$, see \cite[Cor.~12.18]{Maggi2012}
or \cite[Thm.~4.9]{HenrotBook}), and weakly lower semicontinuous; the
direct method gives existence and uniqueness of $u_\Om$.
\end{remark}

\begin{definition}[Torsional rigidity and Saint--Venant energy]
\label{def:Trigidity}
The \emph{torsional rigidity} of $\Om$ is
\begin{equation}\label{eq:Tdef}
   T(\Om):=\int_{\Om}u_\Om\,dx.
\end{equation}
The \emph{BDV-normalised Saint--Venant energy} of $\Om$ is
\begin{equation}\label{eq:Edef}
   E(\Om):=\min_{v\in H^{1}_{0}(\Om)}J(v)
   =\tfrac12\!\int_{\Om}\!|\nabla u_\Om|^{2}\,dx-\!\int_{\Om}\!u_\Om\,dx.
\end{equation}
\end{definition}

\begin{lemma}[Energy--rigidity identity]\label{lem:E-T}
For every open $\Om\subset\R^n$ with $0<|\Om|<\infty$,
\begin{equation}\label{eq:E-T}
   E(\Om)=-\tfrac12\,T(\Om).
\end{equation}
In particular $E(\Om)\le 0$ on every admissible set, and $T(\Om)>0$
unless $|\Om|=0$.
\end{lemma}

\begin{proof}
Test the weak equation $-\Delta u_\Om=1$ against $u_\Om\in H^{1}_{0}(\Om)$
to get $\int_\Om|\nabla u_\Om|^2\,dx=\int_\Om u_\Om\,dx=T(\Om)$.
Substituting into~\eqref{eq:Edef} gives
$E(\Om)=\tfrac12 T(\Om)-T(\Om)=-\tfrac12 T(\Om)$.
\end{proof}

\begin{remark}[Two coexisting conventions]\label{rem:T-vs-E}
Some sources (e.g.\ \cite[\S1]{Talenti1976}, classical analysis
literature) call $T(\Om)$ the ``torsion energy''. The BDV chain
(\cite{BDV}; \texttt{Final/SaintVenantAssembly.tex}; the global
reduction in \texttt{Final-BoundedReduction}; the Kohler--Jobin
transfer in \texttt{Final-KohlerJobinTransfer}) and the Plan~2
WIP files (\texttt{WIP\_LevelSetIdentity.tex},
\texttt{WIP\_master.tex}) reserve ``Saint--Venant energy'' for the
BDV-normalised functional $E=-T/2$, which is the minimised value of
the variational functional $J$. We adopt the second convention
throughout: $T(\Om)$ is the torsional \emph{rigidity}; $E(\Om)$ is
the BDV-normalised Saint--Venant \emph{energy}. The identity
\eqref{eq:E-T} converts between the two.
\end{remark}

\begin{lemma}[Scaling and ball values]\label{lem:Tscaling}
For every $t>0$,
\begin{equation}\label{eq:Tscaling}
   T(t\Om)=t^{n+2}T(\Om),\qquad E(t\Om)=t^{n+2}E(\Om).
\end{equation}
For the unit ball $B_1$, the explicit torsion function is
$u_{B_1}(x)=(1-|x|^2)/(2n)$, hence
\begin{equation}\label{eq:Tn}
   T_n:=T(B_1)=\int_{B_1}\frac{1-|x|^2}{2n}\,dx=\frac{\omega_n}{n(n+2)},
   \qquad E(B_1)=-\frac{\omega_n}{2n(n+2)}.
\end{equation}
\end{lemma}

\begin{proof}
Scaling: with $v_t(x):=t^2 u_\Om(x/t)\in H^1_0(t\Om)$, one verifies
$-\Delta v_t=1$ in $t\Om$, so by uniqueness $u_{t\Om}=v_t$; then
$T(t\Om)=\int_{t\Om}v_t=t^{n+2}\int_\Om u_\Om=t^{n+2}T(\Om)$, and
\eqref{eq:E-T} extends to $E$. For the ball formula, a direct
computation gives $\Delta((1-|x|^2)/(2n))=-1$ and the boundary
value $0$. Integrating in radial coordinates,
$\int_{B_1}(1-|x|^2)/(2n)\,dx
=(1/(2n))[\omega_n-n\omega_n/(n+2)]=\omega_n/(n(n+2))$.
\end{proof}

\subsubsection{Saint--Venant inequality and Saint--Venant deficit}
\label{ssec:SV-deficit}

The classical Saint--Venant inequality
\cite[Thm.~3.3]{PolyaSzego1951} states that the ball maximises torsional
rigidity at fixed volume.

\begin{theorem}[Saint--Venant inequality, classical]\label{thm:SVclass}
For every open $\Om\subset\R^n$ with $0<|\Om|<\infty$, and with
$\Bstar$ the open ball centred at the origin with $|\Bstar|=|\Om|$,
\begin{equation}\label{eq:SVclass}
   T(\Om)\le T(\Bstar),
   \qquad\text{equivalently}\qquad
   E(\Om)\ge E(\Bstar).
\end{equation}
Equality holds iff $\Om$ is (a.e.\ equal to) a ball.
\end{theorem}

\begin{proof}[Citation]
\cite[Thm.~1.1]{PolyaSzego1951}, also obtained by Schwarz symmetrisation
combined with the P\'olya--Szeg\H{o} principle, or as a corollary of
Talenti's rearrangement comparison (Theorem~\ref{thm:talenti}
below): on the unit-volume normalisation, the latter gives
$u^*_\Om\le u^*_{\Bstar}$ pointwise on $[0,|\Om|]$, hence
$T(\Om)=\int_0^{|\Om|}u^*_\Om(s)ds\le\int_0^{|\Om|}u^*_{\Bstar}(s)ds=T(\Bstar)$.
\end{proof}

\begin{definition}[Saint--Venant deficit $\deltaT$]\label{def:deltaT}
For open $\Om\subset\R^n$ with $0<|\Om|<\infty$ and $\Bstar$ the
ball with $|\Bstar|=|\Om|$, the (BDV-normalised) Saint--Venant
deficit is
\begin{equation}\label{eq:deltaT-def}
   \deltaT(\Om):=E(\Om)-E(\Bstar)\ge 0,
\end{equation}
equivalently $\deltaT(\Om)=\tfrac12(T(\Bstar)-T(\Om))$.
\end{definition}

\begin{remark}[Scale-invariant form]\label{rem:Dform}
The scale-invariant version is
\begin{equation}\label{eq:Dform}
   D(\Om):=E(\Om)\,|\Om|^{-(n+2)/n}-E(B_1)\,|B_1|^{-(n+2)/n}\ge 0,
\end{equation}
which for $|\Om|=\omega_n$ reduces to
$D(\Om)=\omega_n^{-(n+2)/n}\deltaT(\Om)$.
Both formulations are
used: $\deltaT$ in the Plan~2 chain and the bounded Saint--Venant
statements; $D$ in the bounded reduction (Plan~1, \S\,IV).
\end{remark}

\subsubsection{Fraenkel asymmetry}\label{ssec:Fraenkel}

\begin{definition}[Fraenkel asymmetry]\label{def:Fraenkel}
For open $\Om\subset\R^n$ with $0<|\Om|<\infty$, the
\emph{Fraenkel asymmetry} is
\begin{equation}\label{eq:Fraenkel-def}
   \Fr(\Om):=\inf_{x_0\in\R^n}
        \frac{|\Om\,\triangle\,(\Bstar+x_0)|}{|\Bstar|},
\end{equation}
where $\Bstar$ is the ball centred at the origin with $|\Bstar|=|\Om|$.
\end{definition}

\begin{lemma}[Range of $\Fr$]\label{lem:Fraenkel-range}
$\Fr(\Om)\in[0,2)$, with $\Fr(\Om)=0$ iff $\Om$ is (a.e.\ equal to) a
translate of $\Bstar$.
\end{lemma}

\begin{proof}
Non-negativity is immediate. The strict bound $\Fr<2$ follows from
$|\Om\,\triangle\,(\Bstar+x_0)|=|\Om|+|\Bstar|-2|\Om\cap(\Bstar+x_0)|$
and the existence (by translation continuity of Lebesgue measure)
of $x_0\in\R^n$ with $|\Om\cap(\Bstar+x_0)|>0$, so
$|\Om\,\triangle\,(\Bstar+x_0)|<2|\Bstar|$. The infimum is attained
because the map $x_0\mapsto|\Om\,\triangle\,(\Bstar+x_0)|$ is
continuous and tends to $2|\Bstar|$ as $|x_0|\to\infty$
(\cite[Lem.~3.2]{Maggi2008}).
\end{proof}

\begin{remark}[Translation invariance]\label{rem:trans}
Both $\deltaT$, $\deltaFK$, and $\Fr$ are translation invariant; all
estimates below are stated up to translation.
\end{remark}

\subsubsection{Faber--Krahn deficit and eigenvalue normalisation}
\label{ssec:Faber-Krahn}

\begin{definition}[First Dirichlet eigenvalue]\label{def:lambda1}
For open $\Om\subset\R^n$ with $0<|\Om|<\infty$,
\begin{equation}\label{eq:lambda1-def}
   \lambda_1(\Om):=\min_{\substack{v\in H^{1}_{0}(\Om)\\v\not\equiv 0}}
   \frac{\int_\Om|\nabla v|^2\,dx}{\int_\Om v^2\,dx}.
\end{equation}
The minimum is attained (by direct method on
$\{v:\|v\|_{L^2}=1\}$, using the Poincar\'e--Sobolev inequality on
finite-measure sets), and the eigenvalue scales as
$\lambda_1(t\Om)=t^{-2}\lambda_1(\Om)$ for $t>0$. We write
$L_n:=\lambda_1(B_1)$ for the unit-ball first eigenvalue. Numerically,
$L_n=j_{n/2-1,1}^2$, the square of the first positive zero of the
Bessel function $J_{n/2-1}$ \cite[\S1.3]{HenrotBook}.
\end{definition}

\begin{theorem}[Faber--Krahn inequality]\label{thm:FK}
For every open $\Om\subset\R^n$ with $0<|\Om|<\infty$,
\begin{equation}\label{eq:FK}
   \lambda_1(\Om)\ge\lambda_1(\Bstar)
   =L_n\,\bigl(|\Bstar|/\omega_n\bigr)^{-2/n},
\end{equation}
with equality iff $\Om$ is (a.e.\ equal to) a ball.
\end{theorem}

\begin{proof}[Citation]
The classical Faber--Krahn theorem \cite{Faber1923,Krahn1925}, proved
by Schwarz symmetrisation and the P\'olya--Szeg\H{o} inequality. See
\cite[Thm.~3.2.1]{HenrotBook}.
\end{proof}

\begin{definition}[Faber--Krahn deficit $\deltaFK$]
\label{def:deltaFK}
For open $\Om\subset\R^n$ with $0<|\Om|<\infty$ and $\Bstar$ the ball
with $|\Bstar|=|\Om|$, the scale-normalised Faber--Krahn deficit is
\begin{equation}\label{eq:deltaFK-def}
   \deltaFK(\Om)
   :=\Bigl(\frac{|\Om|}{\omega_n}\Bigr)^{2/n}
     \bigl(\lambda_1(\Om)-\lambda_1(\Bstar)\bigr)\ge 0.
\end{equation}
\end{definition}

\begin{lemma}[Scaling invariance of $\deltaFK$]\label{lem:deltaFK-scale}
$\deltaFK(t\Om)=\deltaFK(\Om)$ for every $t>0$.
\end{lemma}

\begin{proof}
$\lambda_1(t\Om)=t^{-2}\lambda_1(\Om)$ and
$|t\Om|^{2/n}=t^2|\Om|^{2/n}$ give
$|t\Om|^{2/n}\lambda_1(t\Om)=|\Om|^{2/n}\lambda_1(\Om)$. Similarly for
$\Bstar$.
\end{proof}

\begin{remark}[Why this normalisation]\label{rem:deltaFK-form}
Both Plan~1 and Plan~2 produce, after Kohler--Jobin transfer, a sharp
quantitative inequality of the form
$\deltaFK(\Om)\ge c_{\mathrm{FK}}(n)\,\Fr(\Om)^2$. The factor
$(|\Om|/\omega_n)^{2/n}$ in~\eqref{eq:deltaFK-def} makes both sides
scale-invariant, mirroring~\eqref{eq:Dform}; see
\texttt{Final-KohlerJobinTransfer}, eq.\,(eq:fkout).
\end{remark}

\subsubsection{Talenti pointwise comparison}\label{ssec:Talenti}

The fundamental rearrangement comparison for the torsion function is
the Talenti pointwise inequality \cite{Talenti1976}. We write
$u^*\colon[0,|\Om|]\to[0,\|u\|_\infty]$ for the decreasing
rearrangement of $u=u_\Om$, $u^*(s):=\inf\{t\ge 0\,:\,m(t)\le s\}$,
where $m(t):=|\{u>t\}|$.

\begin{theorem}[Talenti, 1976]\label{thm:talenti}
Let $\Om\subset\R^n$ be open with $0<|\Om|<\infty$, and let
$u_\Om\in H^{1}_{0}(\Om)$ be the torsion function of
Definition~\ref{def:utorsion}. Let $u_{\Bstar}$ be the torsion
function of the ball $\Bstar$ with $|\Bstar|=|\Om|$. Then
\begin{equation}\label{eq:talenti}
   u^*_\Om(s)\le u^*_{\Bstar}(s)
   =\frac{R_\Om^2-(s/\omega_n)^{2/n}}{2n},
   \qquad s\in[0,|\Om|],
\end{equation}
where $R_\Om:=(|\Om|/\omega_n)^{1/n}$ is the radius of $\Bstar$. In
particular,
\begin{equation}\label{eq:talenti-sup}
   \|u_\Om\|_{L^\infty(\Om)}\le u^*_{\Bstar}(0)=\frac{R_\Om^2}{2n}.
\end{equation}
\end{theorem}

\begin{proof}[Citation]
\cite[Thm.~1]{Talenti1976}. The argument uses isoperimetric
rearrangement of super-level sets and the coarea formula; it is
constructive and gives equality iff $u_\Om$ is (translate of) the
ball torsion function, hence iff $\Om$ is a ball.
\end{proof}

\begin{remark}[Consequence: $L^\infty$ bound used in Plan~2]
\label{rem:talenti-Linf}
The sup-norm bound~\eqref{eq:talenti-sup} is the only reason that
the level-set deficit identity
(\texttt{WIP\_LevelSetIdentity.tex}, Thm.\,1) holds for arbitrary
open sets of finite measure without an explicit boundedness
hypothesis. Throughout Plan~2 we use the normalised volume
$|\Om|=\omega_n$, so $R_\Om=1$ and
$\|u_\Om\|_\infty\le 1/(2n)$.
\end{remark}

\subsubsection{Volume--radius foliation \texorpdfstring{$\Erho=\{u>t(\rho)\}$}{E\_rho=\{u>t(rho)\}}}
\label{ssec:foliation}

We now record the volume-radius parametrisation of the super-level
sets of $u_\Om$, used throughout Plan~2.

\begin{definition}[Distribution function and its inverse]
\label{def:t-rho}
For $\Om$ open with $|\Om|=\omega_n$ (the Plan~2 normalisation),
$E_t:=\{u_\Om>t\}$ and $m(t):=|E_t|$. By the BV coarea formula and
Lemma~\ref{lem:coarea-mprime} below, $m$ is right-continuous and
non-increasing on $[0,\|u\|_\infty]$ with $m(0^+)\le\omega_n$ and
$m(\|u\|_\infty)=0$. Define
\begin{equation}\label{eq:trho}
   t(\rho):=\inf\bigl\{t\ge 0\,:\,m(t)\le\omega_n\rho^n\bigr\},
   \qquad\rho\in[0,1],
\end{equation}
and the volume-radius foliation
\begin{equation}\label{eq:Erho}
   \Erho:=\{u_\Om>t(\rho)\},\qquad
   |\Erho|=\omega_n\rho^n,\quad
   \mu(d\rho):=(-t'(\rho))\,d\rho.
\end{equation}
\end{definition}

\begin{lemma}[Properties of $t(\rho)$]\label{lem:trho-props}
Let $\Om$ open with $|\Om|=\omega_n$ and $u=u_\Om$. Then
\begin{enumerate}
\item[(i)] $\rho\mapsto t(\rho)$ is non-increasing on $[0,1]$, with
$t(1)=0$ and $0\le t(\rho)\le\|u\|_\infty\le 1/(2n)$ by
\eqref{eq:talenti-sup}.
\item[(ii)] $t$ is absolutely continuous on every interval
$[\rstar,1]\subset(0,1]$.
\item[(iii)] On a co-null set in $[0,1]$,
$t'(\rho)\le 0$, and
\begin{equation}\label{eq:dt-bound}
   0\le -t'(\rho)\le\frac{1}{n}
   \qquad\text{for a.e.\ }\rho\in[0,1].
\end{equation}
\end{enumerate}
\end{lemma}

\begin{proof}
(i) and the boundary value $t(1)=0$ follow from $m(0^+)\le\omega_n$
and $m(t(\rho))\le\omega_n\rho^n$; the sup-norm bound is
Theorem~\ref{thm:talenti}.

(ii) The map $m$ is, by Lemma~\ref{lem:coarea-mprime}, absolutely
continuous on $[\eps,\|u\|_\infty]$ for every $\eps>0$ (its derivative
$-\alpha(t)$ is in $L^1_{\mathrm{loc}}$). The inverse of a strictly
decreasing absolutely continuous function with non-vanishing
derivative on an open set is absolutely continuous; on plateau
levels (where $m$ is constant) $t$ has a jump of measure zero and
the equation $|E_\rho|=\omega_n\rho^n$ pins $t(\rho)$ on the
right-continuous representative. On the foliation interval
$[\rstar,1]$ this gives absolute continuity of $t$.

(iii) Differentiating $m(t(\rho))=\omega_n\rho^n$ at a Lebesgue point
of $\rho\mapsto t'(\rho)$ at which $m$ is differentiable at
$t(\rho)$, by the chain rule for absolutely continuous functions and
Lemma~\ref{lem:coarea-mprime},
$m'(t(\rho))\,t'(\rho)=n\omega_n\rho^{n-1}$, i.e.\
\begin{equation}\label{eq:tprime-formula}
   -t'(\rho)=\frac{n\omega_n\rho^{n-1}}{\alpha(t(\rho))},
   \qquad\alpha(t):=\!\int_{\partial^*E_t}\!|\nabla u|^{-1}\,d\mathcal H^{n-1}.
\end{equation}
By the Cauchy--Schwarz inequality on $\partial^*E_t$ and the
isoperimetric inequality $P(E_t)\ge n\omega_n^{1/n}m(t)^{1-1/n}$,
\begin{equation}\label{eq:alpha-LB}
   \alpha(t)\cdot\beta(t)\ge P(E_t)^2\ge n^2\omega_n^{2/n}m(t)^{2-2/n},
   \qquad
   \beta(t):=\!\int_{\partial^*E_t}\!|\nabla u|\,d\mathcal H^{n-1}=m(t),
\end{equation}
the equality $\beta(t)=m(t)$ being the flux identity of
Lemma~\ref{lem:flux-id} below. Combined with $m(t(\rho))=\omega_n\rho^n$,
$\alpha(t(\rho))\ge n^2\omega_n^{2/n}(\omega_n\rho^n)^{2-2/n}/m(t(\rho))
=n^2\omega_n\rho^{n-2}\cdot(\omega_n\rho^n)^{1-2/n}/(\omega_n\rho^n)
=n\omega_n\rho^{n-1}\cdot n$, hence
$-t'(\rho)\le n\omega_n\rho^{n-1}/(n^2\omega_n\rho^{n-1})=1/n$.
\end{proof}

\begin{remark}[Where $-t'\le 1/n$ is used]
\label{rem:H3}
The bound \eqref{eq:dt-bound} is the hypothesis (H3) of the abstract
weighted metric trace lemma (\texttt{WIP\_WeightedMetricTrace.tex},
Thm.\,5.1): it bounds the density of $\mu$ uniformly.
\end{remark}

The coarea and flux identities used in the proof of
Lemma~\ref{lem:trho-props}(iii) are recorded for reference (proved in
the geometric measure theory preliminaries,
\S\ref{sec:prelim-gmt}):

\begin{lemma}[Coarea form of $-m'$]\label{lem:coarea-mprime}
For a.e.\ $t\in(0,\|u\|_\infty)$, $-m'(t)=\alpha(t)$ as defined in
\eqref{eq:tprime-formula}. $m$ is absolutely continuous on
$[\eps,\|u\|_\infty]$ for every $\eps>0$.
\end{lemma}

\begin{proof}[Citation]
\cite[Thm.~13.1]{Maggi2012}, \cite[Thm.~3.40]{AFP2000}.
\end{proof}

\begin{lemma}[Flux--volume identity]\label{lem:flux-id}
For a.e.\ $t\in(0,\|u\|_\infty)$, $E_t$ has finite perimeter and
$\beta(t)=m(t)$.
\end{lemma}

\begin{proof}[Citation]
By Lipschitz truncation of the weak equation $-\Delta u=1$ against the
cutoff $\phi_\eps(u-t)\in H^{1}_{0}(\Om)$ and the finite-perimeter
divergence theorem; cf.\ \cite[Thm.~16.3]{Maggi2012}.
\end{proof}

\subsubsection{Profile-gap estimate used by Plan~2}\label{ssec:profile-gap}

We now record the profile-gap estimate that converts an
$L^\infty$-control of $u^*_\Om-u^*_{\Bstar}$ on a near-boundary slab
into a positive lower bound for $|T_\eta|=v(L-2\eta)-v(L-\eta)$, hence
into the hypothesis~(H2) of the abstract weighted metric trace lemma.

We work in the Plan~2 normalisation: $|\Om|=\omega_n=:L$,
$R_\Om=R=1$, and we write $b(s):=u^*_{B_1}(s)$, $v(s):=u^*_\Om(s)$,
$h(s):=b(s)-v(s)$.

\begin{lemma}[Profile gap; integrated form]\label{lem:profile-gap-int}
Let $\Om\subset\R^n$ be open with $|\Om|=\omega_n$. Then for every
$s\in(0,\omega_n]$,
\begin{equation}\label{eq:profile-gap-int}
   0\le h(s)=b(s)-v(s)\le\frac{2\deltaT(\Om)}{s}.
\end{equation}
In particular, for every $\theta\in(0,1)$,
\begin{equation}\label{eq:profile-gap-slab}
   \|h\|_{L^\infty([\theta\omega_n,\omega_n])}
   \le\frac{2\deltaT(\Om)}{\theta\omega_n}.
\end{equation}
\end{lemma}

\begin{proof}
By Talenti's pointwise comparison (Theorem~\ref{thm:talenti}),
$v(s)\le b(s)$, so $h(s)\ge 0$ and $h$ is the difference of
non-increasing absolutely continuous functions, hence absolutely
continuous on $(0,\omega_n]$ with $h(0^+)\le b(0)<\infty$ and
$h(\omega_n)=0$.

For the upper bound, integrate Talenti's level-set differential
inequality (cf.~\cite[\S4]{Talenti1976},
\texttt{Plan 2/level-set-deficit-identity.md} §4):
$v'(s)\ge b'(s)$ for a.e.\ $s$, equivalently $h'(s)\le 0$. Since
$h\ge 0$ is non-increasing,
$h(s)\cdot s\le\int_0^s h(\sigma)\,d\sigma\le\int_0^{\omega_n}h$. By
equimeasurability,
\[
   \int_0^{\omega_n}h(\sigma)\,d\sigma
   =\int_0^{\omega_n}b(\sigma)\,d\sigma-\int_0^{\omega_n}v(\sigma)\,d\sigma
   =T(B_1)-T(\Om)=2\deltaT(\Om),
\]
using $|\Om|=\omega_n=|B_1|$ and Lemma~\ref{lem:E-T}. Hence
$h(s)\le 2\deltaT(\Om)/s$.
\end{proof}

\begin{lemma}[Foliation length on a boundary slab]
\label{lem:slab-length}
Let $\Om\subset\R^n$ open with $|\Om|=\omega_n$. For
$\eta\in(0,\omega_n/4)$ define the volume slab $S_\eta:=[\omega_n-2\eta,\omega_n-\eta]$
and the level slab
$T_\eta:=v(S_\eta)=[v(\omega_n-\eta),v(\omega_n-2\eta)]$. There exist
constants $c_n,C_n>0$ depending only on $n$ such that
\begin{equation}\label{eq:slab-b-bounds}
   c_n\eta\le b(s)\le C_n\eta\qquad\text{for every }s\in S_\eta,
\end{equation}
and if $\deltaT(\Om)\le c_n\eta$ then
\begin{equation}\label{eq:slab-T-eta}
   |T_\eta|=v(\omega_n-2\eta)-v(\omega_n-\eta)\ge c_n\eta.
\end{equation}
The constants are explicit: $c_n=1/(C\cdot n)$, $C_n=C\cdot n^{-1}$,
$C$ a universal dimensional constant.
\end{lemma}

\begin{proof}
Set $R=1$, $L=\omega_n$. From the explicit ball profile
$b(s)=(1-(s/\omega_n)^{2/n})/(2n)$, expanding
$1-(1-\eta/\omega_n)^{2/n}$ to first order in $\eta$ for $\eta\le\omega_n/4$
gives
$c_n\eta/\omega_n\le 1-(1-\eta/\omega_n)^{2/n}\le C_n\eta/\omega_n$
with $c_n=1/n$, $C_n=4/n$ (using
$1-(1-x)^{2/n}\ge(2/n)x(1-x)\ge x/n$ for $x\in[0,1/2]$, and
$1-(1-x)^{2/n}\le(2/n)x\cdot(1-x)^{2/n-1}\le 4x/n$ for $n\ge 2$).
Multiplying by $1/(2n)$ gives \eqref{eq:slab-b-bounds}.

For \eqref{eq:slab-T-eta}, write
$v(\omega_n-2\eta)-v(\omega_n-\eta)
=[b(\omega_n-2\eta)-b(\omega_n-\eta)]
 -[h(\omega_n-2\eta)-h(\omega_n-\eta)]$.
The first bracket is, by the explicit ball profile and
$\omega_n-2\eta\ge\omega_n/2$, at least
$c_n\eta$ for a dimensional $c_n$. The second bracket is bounded by
$2\|h\|_{L^\infty([\omega_n/2,\omega_n])}\le 4\deltaT(\Om)/\omega_n$
via \eqref{eq:profile-gap-slab}. The smallness hypothesis
$\deltaT(\Om)\le c_n\eta$ (after halving $c_n$) ensures the first
bracket dominates the second.
\end{proof}

\begin{remark}[Adaptation source]
Lemma~\ref{lem:profile-gap-int} is the integrated form of the
Talenti-comparison consequence
$0\le b(s)-v(s)\le 2\deltaT/s$ derived from
\cite[Thm.~1]{Talenti1976} and the Saint--Venant deficit identity
$\Gamma(\Om)=2(E(\Om)-E(\Bstar))/|\Om|$. Lemma~\ref{lem:slab-length}
provides the $|T_\eta|$-bound used to verify hypothesis~(H2) in the
abstract trace lemma of \S\ref{subsec:weighted-metric-trace}.
\end{remark}

\subsubsection{Bad-set Lebesgue estimate used by Plan~2}\label{ssec:badset}

The bad-set estimate converts the integrated Talenti gap into a
Lebesgue control on the set where the foliation slows down. It is
the input that the weighted metric trace lemma needs to convert the
weighted kinetic bound
$\int|\dot F_\rho|_\X^2\,d\mu\le C\deltaT$ into the unweighted kinetic
bound $\int|\dot F_\rho|_\X^2\,d\rho\le C'\deltaT$
(cf.\ \S\ref{subsec:weighted-metric-trace}).

\begin{lemma}[Integrated Talenti gap on $-t'$]\label{lem:talenti-gap}
Let $\Om\subset\R^n$ be open with $|\Om|=\omega_n$ and
$\rstar\in(0,1)$. Let $t_B$ denote the inverse distribution
function of $u_{B_1}$,
$t_B(\rho)=u^*_{B_1}(\omega_n\rho^n)=(1-\rho^2)/(2n)$, so
$-t_B'(\rho)=\rho/n$. There is $K_1=K_1(n,\rstar)>0$ such that
\begin{equation}\label{eq:talenti-gap-int}
   \int_{\rstar}^{1}\bigl|(-t_B'(\rho))-(-t'(\rho))\bigr|\,d\rho
   \le K_1\,\deltaT(\Om).
\end{equation}
\end{lemma}

\begin{proof}
Since both $t,t_B$ are non-increasing with $t(1)=t_B(1)=0$,
$\psi(\rho):=(-t_B'(\rho))-(-t'(\rho))$ satisfies
$\int_\rho^1\psi(r)\,dr=t(\rho)-t_B(\rho)=v(\omega_n\rho^n)-b(\omega_n\rho^n)
=-h(\omega_n\rho^n)\le 0$ for every $\rho\in[\rstar,1]$
(by Talenti, $v\le b$, hence $h\ge 0$). However $\psi$ need not
have a sign; the proof of Lemma~\ref{lem:profile-gap-int} shows
$h\ge 0$ globally, so the function $-h(\omega_n\rho^n)$ is
$\le 0$ on $[\rstar,1]$.

Substitute $s=\omega_n\rho^n$, $ds=n\omega_n\rho^{n-1}d\rho$:
\[
   \int_{\rstar}^1\psi(\rho)\,d\rho
   =\bigl[\,t(\rho)-t_B(\rho)\,\bigr]_{\rstar}^1
   =-\bigl(t(\rstar)-t_B(\rstar)\bigr)
   =h(\omega_n\rstar^n)
   \le \frac{2\deltaT(\Om)}{\omega_n\rstar^n}
\]
by Lemma~\ref{lem:profile-gap-int}. Since $\psi$ need not be signed,
this controls only $\int\psi$ not $\int|\psi|$. We restore the
absolute value by a second comparison. On any subinterval
$[\rho_1,\rho_2]\subset[\rstar,1]$,
\[
   \Bigl|\!\int_{\rho_1}^{\rho_2}\!\psi(r)\,dr\Bigr|
   =\bigl|(t(\rho_1)-t(\rho_2))-(t_B(\rho_1)-t_B(\rho_2))\bigr|
   =|h(\omega_n\rho_1^n)-h(\omega_n\rho_2^n)|
   \le h(\omega_n\rstar^n)
\]
because $h\ge 0$ is non-increasing. By the
Hahn--Jordan decomposition of $\psi\,d\rho$ on $[\rstar,1]$ as
$d\psi^+-d\psi^-$,
\[
   \int_{\rstar}^1|\psi|\,d\rho
   =\psi^+([\rstar,1])+\psi^-([\rstar,1])
   \le 2\,\sup_{[\rho_1,\rho_2]\subset[\rstar,1]}
   \Bigl|\int_{\rho_1}^{\rho_2}\psi\Bigr|
   \le 2\,h(\omega_n\rstar^n)
   \le\frac{4\deltaT(\Om)}{\omega_n\rstar^n},
\]
which is \eqref{eq:talenti-gap-int} with
$K_1=4/(\omega_n\rstar^n)$.
\end{proof}

\begin{lemma}[Bad-set Lebesgue estimate; (H3-bad)]
\label{lem:badset}
Let $\Om\subset\R^n$ open with $|\Om|=\omega_n$, $\rstar\in(0,1)$.
Set
\[
   \tau_0:=\frac{\rstar}{2n},
   \qquad
   B_{\tau_0}:=\{\rho\in[\rstar,1]\,:\,-t'(\rho)<\tau_0\}.
\]
Then
\begin{equation}\label{eq:badset}
   |B_{\tau_0}|\le K_2(n,\rstar)\,\deltaT(\Om),
   \qquad K_2(n,\rstar):=\frac{8n}{\omega_n\rstar^{n+1}}.
\end{equation}
\end{lemma}

\begin{proof}
On $B_{\tau_0}$, $-t'(\rho)<\rstar/(2n)$ and $-t_B'(\rho)=\rho/n\ge\rstar/n$
for $\rho\ge\rstar$. Hence
$\psi(\rho)=(-t_B'(\rho))-(-t'(\rho))>\rstar/n-\rstar/(2n)=\rstar/(2n)$
on $B_{\tau_0}$. Integrating and applying
Lemma~\ref{lem:talenti-gap},
\[
   \frac{\rstar}{2n}\,|B_{\tau_0}|
   \le\int_{B_{\tau_0}}\psi(\rho)\,d\rho
   \le\int_{\rstar}^{1}|\psi(\rho)|\,d\rho
   \le K_1\,\deltaT(\Om),
\]
so $|B_{\tau_0}|\le (2n/\rstar)K_1\deltaT(\Om)$. Substituting
$K_1=4/(\omega_n\rstar^n)$ gives \eqref{eq:badset}.
\end{proof}

\begin{remark}[Adaptation source]
Lemma~\ref{lem:badset} is the bad-set lemma used in
\S\ref{subsec:weighted-metric-trace}. Lemma~\ref{lem:talenti-gap} is
the ``integrated Talenti gap'' invoked there; the constants
$K_1(n,\rstar),K_2(n,\rstar)$ are made explicit here. No published
reference for this exact statement was located; the proof is
self-contained from Talenti's comparison and the level-set deficit
identity.
\end{remark}

\subsubsection{Kohler--Jobin inequality}\label{ssec:KJ}

\begin{theorem}[Kohler--Jobin inequality; \protect\cite{KohlerJobin1978,Brasco2014}]
\label{thm:KJ}
For every open bounded $\Om\subset\R^n$ with $\lambda_1(\Om)<\infty$,
\begin{equation}\label{eq:KJ}
   T(\Om)\cdot\lambda_1(\Om)^{(n+2)/2}
   \ge T(B_1)\cdot\lambda_1(B_1)^{(n+2)/2}
   =T_n\,L_n^{(n+2)/2}.
\end{equation}
Equality holds iff $\Om$ is a ball (up to translation and a null set).
\end{theorem}

\begin{proof}[Citation]
\cite[Thm.~3.3]{Brasco2014}, giving a modern proof by Schwarz
symmetrisation; original \cite{KohlerJobin1978}. The functional
$\Psi(\Om):=T(\Om)\lambda_1(\Om)^{(n+2)/2}$ is scale-invariant
($T(t\Om)=t^{n+2}T(\Om)$ and $\lambda_1(t\Om)=t^{-2}\lambda_1(\Om)$,
so $\Psi(t\Om)=\Psi(\Om)$), hence the right side is the universal
constant $T_n L_n^{(n+2)/2}$ depending only on~$n$.
\end{proof}

\begin{remark}[Normalisation used in the final transfer]
\label{rem:KJ-norm}
The form~\eqref{eq:KJ} is the one used in the Kohler--Jobin transfer
of the present manuscript
(\S\ref{sec:plan1-kj-transfer}). At
fixed volume $|\Om|=|\Bstar|$, rearranging~\eqref{eq:KJ} yields the
pointwise transfer
\begin{equation}\label{eq:KJ-ptwise}
   \lambda_1(\Om)
   \ge\lambda_1(\Bstar)
   \Bigl(\frac{T(\Bstar)}{T(\Om)}\Bigr)^{2/(n+2)},
\end{equation}
equivalently
\begin{equation}\label{eq:KJ-deficit}
   \lambda_1(\Om)-\lambda_1(\Bstar)
   \ge\lambda_1(\Bstar)\Bigl[(1-s)^{-1/\beta}-1\Bigr],
   \qquad
   s:=\frac{T(\Bstar)-T(\Om)}{T(\Bstar)}\in[0,1),
   \quad\beta=\tfrac{n+2}{2}.
\end{equation}
\end{remark}

\begin{lemma}[Elementary one-variable inequality]\label{lem:onevar}
For $p\in(0,1]$ and $s\in[0,1)$,
\begin{equation}\label{eq:onevar}
   (1-s)^{-p}-1\ge ps.
\end{equation}
\end{lemma}

\begin{proof}
$f(s):=(1-s)^{-p}-1$ satisfies $f(0)=0$ and
$f'(s)=p(1-s)^{-(p+1)}\ge p$ for $s\in[0,1)$. By the fundamental
theorem of calculus, $f(s)=\int_0^s f'(\tau)\,d\tau\ge ps$.
\end{proof}

Combining~\eqref{eq:KJ-deficit} with Lemma~\ref{lem:onevar} for
$p=1/\beta=2/(n+2)$ gives, in the small-deficit regime
$s\le 1/2$, equivalently $T(\Bstar)-T(\Om)\le T(\Bstar)/2$,
equivalently $\deltaT(\Om)\le T(\Bstar)/4$:
\begin{equation}\label{eq:KJ-linear}
   \lambda_1(\Om)-\lambda_1(\Bstar)
   \ge\frac{2\lambda_1(\Bstar)}{(n+2)T(\Bstar)}
       \bigl(T(\Bstar)-T(\Om)\bigr)
   =\frac{4\lambda_1(\Bstar)}{(n+2)T(\Bstar)}\,\deltaT(\Om).
\end{equation}

\begin{remark}[Dimension dependence]\label{rem:KJ-dim}
The Kohler--Jobin multiplier in~\eqref{eq:KJ-linear} depends only on
$n$: by scale invariance,
$2\lambda_1(\Bstar)/((n+2)T(\Bstar))=2L_n/((n+2)T_n)=2nL_n/\omega_n$,
using \eqref{eq:Tn}. This is the universal dimension-only constant
that converts $\deltaT$ to $\deltaFK$ in the final theorem.
\end{remark}

\subsubsection{Summary of constants and constants used downstream}
\label{ssec:prelim-summary}

We record the universal ball constants and the dimension-tracked
constants introduced in this section, for use in the Plan~1 and
Plan~2 chains.

\begin{center}
\renewcommand{\arraystretch}{1.25}
\begin{tabular}{|l|l|l|}\hline
Symbol & Value / dependence & Reference\\\hline
$\omega_n=|B_1|$ & dim.~only & --- \\
$T_n=T(B_1)=\omega_n/(n(n+2))$ & dim.~only & Lemma~\ref{lem:Tscaling}\\
$L_n=\lambda_1(B_1)=j_{n/2-1,1}^2$ & dim.~only & Def.~\ref{def:lambda1}\\
$E(B_1)=-T_n/2=-\omega_n/(2n(n+2))$ & dim.~only & \eqref{eq:E-T}\\
$2L_n/((n+2)T_n)=2nL_n/\omega_n$ & Kohler--Jobin multiplier ($n$ only) & Rmk.~\ref{rem:KJ-dim}\\\hline
$K_1(n,\rstar)=4/(\omega_n\rstar^n)$ & integrated Talenti gap constant & Lem.~\ref{lem:talenti-gap}\\
$K_2(n,\rstar)=8n/(\omega_n\rstar^{n+1})$ & bad-set Lebesgue constant & Lem.~\ref{lem:badset}\\
$\tau_0=\rstar/(2n)$ & bad-set threshold & Lem.~\ref{lem:badset}\\\hline
$c_n,C_n$ & ball-profile constants in slab estimates & Lem.~\ref{lem:slab-length}\\\hline
\end{tabular}
\end{center}

% =====================================================================
% Local bibliography for this section. Agent 17 will merge this into
% the master bibliography (Manuscript/04_references_audit.md).
% DISABLED by Agent 19 (Final Assembly): the master bibliography is in
% Manuscript/references.bib, loaded from main.tex.
% =====================================================================
\iffalse
\begin{thebibliography}{99}

\bibitem{AFP2000}
L.~Ambrosio, N.~Fusco, D.~Pallara,
\emph{Functions of Bounded Variation and Free Discontinuity
Problems}, Oxford University Press, 2000.

\bibitem{BDV}
L.~Brasco, G.~De~Philippis, B.~Velichkov,
\emph{Faber--Krahn inequalities in sharp quantitative form},
Duke Math.\ J.\ \textbf{164} (2015), 1777--1831.

\bibitem{Brasco2014}
L.~Brasco,
\emph{On torsional rigidity and principal frequencies: an invitation
to the Kohler--Jobin rearrangement technique}, ESAIM Control Optim.\
Calc.\ Var.\ \textbf{20} (2014), 315--338.

\bibitem{Faber1923}
G.~Faber, \emph{Beweis, dass unter allen homogenen Membranen von
gleicher Fl\"ache und gleicher Spannung die kreisf\"ormige den
tiefsten Grundton gibt}, Sitzungsber.\ Bayer.\ Akad.\ Wiss.\
M\"unchen, math.-phys.\ Kl.\ (1923), 169--172.

\bibitem{HenrotBook}
A.~Henrot, \emph{Extremum Problems for Eigenvalues of Elliptic
Operators}, Birkh\"auser, 2006.

\bibitem{KohlerJobin1978}
M.-Th.~Kohler-Jobin,
\emph{Une m\'ethode de comparaison isop\'erim\'etrique de fonctionnelles
de domaines de la physique math\'ematique},
J.\ Appl.\ Math.\ Phys.\ (ZAMP) \textbf{29} (1978), 757--766.

\bibitem{Krahn1925}
E.~Krahn, \emph{\"Uber eine von Rayleigh formulierte
Minimaleigenschaft des Kreises}, Math.\ Ann.\ \textbf{94} (1925),
97--100.

\bibitem{Maggi2012}
F.~Maggi, \emph{Sets of Finite Perimeter and Geometric Variational
Problems: An Introduction to Geometric Measure Theory}, Cambridge
University Press, 2012.

\bibitem{Maggi2008}
F.~Maggi, \emph{Some methods for studying stability in
isoperimetric-type problems}, Bull.\ Amer.\ Math.\ Soc.\
\textbf{45} (2008), 367--408.

\bibitem{PolyaSzego1951}
G.~P\'olya, G.~Szeg\H{o}, \emph{Isoperimetric Inequalities in
Mathematical Physics}, Princeton University Press, 1951.

\bibitem{Talenti1976}
G.~Talenti, \emph{Elliptic equations and rearrangements}, Ann.\ Mat.\
Pura Appl.\ (4) \textbf{110} (1976), 353--372.

% Local block references (point to source documents shipped in this
% project; will be replaced by internal cross-references at assembly
% time by Agent~17 / 19).

\bibitem{WIP_LevelSetIdentity}
Building block \textsc{LevelSetIdentity},
\texttt{Plan 2/WIP/WIP\_LevelSetIdentity.tex}.

\bibitem{WIP_WeightedMetricTrace}
Building block \textsc{WeightedMetricTrace},
\texttt{Plan 2/WIP/WIP\_WeightedMetricTrace.tex}.

\bibitem{Plan2-LevelSet}
\emph{Level-set deficit identity for the Nicola--Tilli route},
\texttt{Plan 2/level-set-deficit-identity.md}.

\bibitem{Final-KohlerJobinTransfer}
Building block \textsc{KohlerJobinTransfer},
\texttt{Final-KohlerJobinTransfer}.

\bibitem{Final-BoundedReduction}
Building block \textsc{BoundedReduction},
\texttt{Final-BoundedReduction}.

\bibitem{FigalliMaggi2011}
A.~Figalli, F.~Maggi,
\emph{On the shape of liquid drops and crystals in the small mass
regime}, Arch.\ Rational Mech.\ Anal.\ \textbf{201} (2011),
143--207 (Appendix~A, Theorem~A.1, the Sobolev--Sard property).

\bibitem{DeGiorgi1958}
E.~De~Giorgi, \emph{Su una teoria generale della misura
$(r-1)$-dimensionale in uno spazio ad $r$ dimensioni}, Ann.\ Mat.\
Pura Appl.\ (4) \textbf{36} (1954), 191--213.

\bibitem{Kawohl1985}
B.~Kawohl, \emph{Rearrangements and Convexity of Level Sets in
PDE}, Lecture Notes in Math.\ \textbf{1150}, Springer, 1985.

\end{thebibliography}
\fi
